\documentclass[10pt]{beamer}

\usepackage{lecture-style}
\usepackage{caption}
\usepackage{subcaption}

\title{Biophysics}

\author{\textbf{Jakub Rydzewski}\\[5pt]
  \footnotesize{\it NCU Institute of Physics}}

\date{\footnotesize\textcolor{gray}{Updated: \today}}

\begin{document}

{\setbeamertemplate{footline}{}\frame{\titlepage}}

\begin{frame}{Literature}
\begin{itemize}
\setlength\itemsep{1em}
  \item M.~E. Tuckerman, \textit{Statistical Mechanics: Theory and Simulation}, Oxford University Press (2016).
  \item D. Chandler, \textit{Introduction to Modern Statistical Mechanics}, Oxford University Press (1987).
\end{itemize}
\end{frame}


\begin{frame}{}
  \begin{center}
    \fontsize{25pt}{6}\selectfont\vspace{1.2cm}
    \textcolor{subtitle}{Monte Carlo Methods}
  \end{center}
\end{frame}

\begin{frame}{Monte Carlo Methods}
The main steps of the Monte Carlo sampling are:
\vspace*{0.5cm}
\begin{itemize}
\setlength\itemsep{1em}
  \item Choose randomly a new configuration in phase space based on a Markov chain.
  \item Accept or reject the new configuration, depending on the strategy used.
  \item Compute the physical quantity and add it to the averaging procedure.
  \item Repeat the previous steps until convergence.
\end{itemize}
\end{frame}

\begin{frame}{Monte Carlo Methods}

\texttt{code/importance-sampling.py}

\begin{equation}
  I = \int_0^1\d x \int_0^{\sqrt{1-x^2}}\d y = \frac{\pi}{4}
\end{equation}
\end{frame}

\begin{frame}{Monte Carlo Methods}
\begin{itemize}
\setlength\itemsep{1em}
  \item The integrals that must be evaluated in equilibrium statistical mechanics are generally of the form:
  \begin{equation}
  \label{eq:general_int}
    I = \int\d \bx\; \phi(\bx) p(\bx),
  \end{equation}
  where $\bx$ is an $n$-dimensional vector, $\phi(\bx)$ is an arbitrary function, and $p(\bx)$ is a function satisfying the properties of a probability distribution:
  \begin{equation}
    p(\bx)\geq 0~~~\mathrm{and}~~~\int\d\bx\; p(\bx) = 1.
  \end{equation}

  \item Eq.~\ref{eq:general_int} represents the ensemble average of a physical observable in equilibrium statistial mechanics.
\end{itemize}
\end{frame}

\begin{frame}{Monte Carlo Methods}
\begin{itemize}
\setlength\itemsep{1em}
  \item Let $\bx_1, \bx_2, \dots, \bx_M$ be a set of $M$ $n$-dimensional vectors that are sampled from $p(\bx)$.

  \item The problem of sampling from $p(\bx)$ is a nontrivial one.

  \item For now, let us assume that such an algorithm exists. Then, an estimator:
  \begin{equation}
    \hat{I}_M = \frac{1}{M}\sum_{i=1}^M \phi(\bx_i)
  \end{equation}
  is such that:
  \begin{equation}
  \label{eq:est}
    \lim_{M\rightarrow\infty} \hat{I}_M = I.
  \end{equation}

  \item Eq.~\ref{eq:est} is guaranteed by the \textit{central limit theorem}.\fref{For the derivation, see M.~E. Tuckerman, \textit{Statistical Mechanics: Theory and Simulation}, Oxford University Press (2016), p.~281.}
\end{itemize}
\end{frame}


\begin{frame}{Monte Carlo Methods}
\textbf{Central Limit Theorem}\vspace{0.2cm}
\emphbox{%
If $\bx_1, \bx_2, \dots, \bx_n$ are $n$ random variables drawn from a population with overall mean $\mu$ and finite variance $\sigma^2$, and if $\bar{\bx}_n$ is the sample mean, the limiting form of the distribution, $\lim_{n \rightarrow \infty} \sqrt{n} \left( \frac{\bar{\bx}_n - \mu}{\sigma} \right)$, is the standard normal distribution.}
\end{frame}

\begin{frame}{Monte Carlo Methods}
\begin{itemize}
\setlength\itemsep{1em}
  \item Example of sampling a simple distribution given by:
  \begin{equation}
    p(x) = c\e^{-cx},
  \end{equation}
  on the interval $x\in [0, \infty)$.

  \item To sample from $p(x)$, we need $P(X)$ such that:
  \begin{equation}
    \label{eq:cum}
    P(X) = \int_0^X \d x\; c\e^{-cx} = 1 - \e^{-cX},
  \end{equation}
  and equate Eq.~\ref{eq:cum} to a random number $\xi \in [0,1]$, then we can solve for X, which gives:
  \begin{equation}
    X = -\frac{1}{c}\log(1-\xi).
  \end{equation}

  \item In general, we do not have such simple distribution to sample from in statistical mechanics.
\end{itemize}
\end{frame}

\begin{frame}[fragile]{Monte Carlo Methods: Examples}
\begin{block}{Estimator no 1}
\begin{lstlisting}
def pi_estimator(n_samples):
  est = list()
  n_inside = 0; n_all = 0
  for i in range(n_samples):
    x = np.random.uniform(-1.0, 1.0)
    y = np.random.uniform(-1.0, 1.0)
    dist = np.sqrt(np.power(x, 2) +
                   np.power(y, 2))
    if dist <= 1.0:
      n_inside += 1
    n_all += 1
    est.append(4.0 * np.float(n_inside) /
               np.float(n_all))
  return np.array(est)
\end{lstlisting}
\end{block}
\end{frame}

\begin{frame}[fragile]{Monte Carlo Methods: Examples}
\begin{figure}
  \includegraphics[width=0.8\textwidth]{../code/fig/estimator-1.png}
  \caption{Estimator no 1.}
\end{figure}
\end{frame}

\begin{frame}[fragile]{Monte Carlo Methods: Examples}
\begin{block}{Estimator no 2}
\begin{lstlisting}
def pi_estimator_integral(n_samples):
  est = list()
  e = 0.0
  for i in range(n_samples):
    x = np.random.uniform(0.0, 1.0)
    y = 4.0 * np.sqrt(1.0 - np.power(x, 2))
    e += (y - e) / (np.float(i) + 1.0)
    est.append(e)
  return np.array(est)
\end{lstlisting}
\end{block}
\end{frame}

\begin{frame}[fragile]{Monte Carlo Methods: Examples}
\begin{figure}
  \includegraphics[width=0.8\textwidth]{../code/fig/estimator-2.png}
  \caption{Estimator no 2.}
\end{figure}
\end{frame}

\begin{frame}[fragile]{Monte Carlo Methods: Examples}
\begin{block}{Estimator no 3}
\begin{lstlisting}
def pi_estimator_markov_chain(n_samples, step=0.1):
  est = list()
  n_inside = 0; n_all = 0;
  x = np.random.uniform(-1.0, 1.0)
  y = np.random.uniform(-1.0, 1.0)
  for i in range(n_samples):
    dx = np.random.uniform(-step, step)
    dy = np.random.uniform(-step, step)
    xn = x + dx; yn = y + dy
    if abs(xn) < 1.0 and abs(yn) < 1.0:
      x = xn; y = yn
    dist = np.sqrt(np.power(x, 2) + np.power(y, 2))
    if dist <= 1.0: n_inside += 1
    n_all += 1
    est.append(4.0 * np.float(n_inside) /
               np.float(n_all))
  return np.array(est)
\end{lstlisting}
\end{block}
\end{frame}

\begin{frame}[fragile]{Monte Carlo Methods: Examples}
\begin{figure}
  \includegraphics[width=0.8\textwidth]{../code/fig/estimator-3.png}
  \caption{Estimator no 3.}
\end{figure}
\end{frame}


\begin{frame}{M(RT)$^2$}
\textbf{Markov Chain}\vspace{0.2cm}
\emphbox{
  If the vectors $\bx_1, \dots, \bx_M$ are generated sequentially, and $\bx_{i+1}$ is generated only based on the knowledge of $\bx_i$. the sequence is called a Makov chain.
}
\end{frame}

\begin{frame}{M(RT)$^2$}
\begin{itemize}
\setlength\itemsep{1em}
  \item For physical systems, a Markov chain must satisfy the condition of \textit{detailed balance}, which ensures that the Markov process is microscopically reversible.

  \item M(RT)$^2$ is a rejection method.

  \item With the acceptance probability given by:
  \begin{equation}
    A(\bx_i | \bx_j) = \min\left[1, {f(\bx_i)}/{f(\bx_j)}\right],
  \end{equation}
  where $f(\bx)$ is the probability of the system being at $\bx$.

  \item Sampling in the canonical ensemble:\fref{Reminder: $\br$ are microscopic coordinates.}
  \begin{equation}
    A(\br_i | \br_j) = \min\left[1, \e^{-\beta \left[U(\br_i)-U(\br_j)\right]} \right],
  \end{equation}
  as $f(\br) \propto \e^{-\beta U(\br)}$.
\end{itemize}
\end{frame}

\begin{frame}[fragile]{Particle on Double-Well Potential}
\begin{lstlisting}
python mc-double-well.py --n_samples 10000 --temp 5
\end{lstlisting}
\vspace*{-0.3cm}
\begin{figure}
  \includegraphics[width=\textwidth]{../code/fig/dw-pot.png}
\end{figure}
\end{frame}

\begin{frame}{Particle on Double-Well Potential}
\begin{figure}
  \includegraphics[width=\textwidth]{../code/fig/dw-traj.png}
\end{figure}
\end{frame}

\begin{frame}{Particle on Double-Well Potential}
\begin{figure}
  \includegraphics[width=\textwidth]{../code/fig/dw-hist.png}
\end{figure}
\end{frame}

\begin{frame}[fragile]{Particle on the M\"uller-Brown Potential}
\begin{lstlisting}
python mc-potential.py --n_samples 10000 --temp 1.0
\end{lstlisting}
\vspace*{-0.3cm}
\begin{figure}
  \includegraphics[width=\textwidth]{../code/fig/mc-mb-t-1.png}
\end{figure}
\end{frame}

\begin{frame}[fragile]{Particle on the M\"uller-Brown Potential}
\begin{lstlisting}
python mc-potential.py --n_samples 10000 --temp 2.0
\end{lstlisting}
\vspace*{-0.3cm}
\begin{figure}
  \includegraphics[width=\textwidth]{../code/fig/mc-mb-t-2.png}
\end{figure}
\end{frame}

\begin{frame}[fragile]{Particle on the M\"uller-Brown Potential}
\begin{lstlisting}
python mc-potential.py --n_samples 10000 --temp 3.0
\end{lstlisting}
\vspace*{-0.3cm}
\begin{figure}
  \includegraphics[width=\textwidth]{../code/fig/mc-mb-t-3.png}
\end{figure}
\end{frame}

\begin{frame}[fragile]{Particle on the M\"uller-Brown Potential}
\begin{lstlisting}
python mc-potential.py --n_samples 10000 --temp 4.0
\end{lstlisting}
\vspace*{-0.3cm}
\begin{figure}
  \includegraphics[width=\textwidth]{../code/fig/mc-mb-t-4.png}
\end{figure}
\end{frame}

\begin{frame}[fragile]{Particle on the M\"uller-Brown Potential}
\begin{lstlisting}
python mc-potential.py --n_samples 10000 --temp 5.0
\end{lstlisting}
\vspace*{-0.3cm}
\begin{figure}
  \includegraphics[width=\textwidth]{../code/fig/mc-mb-t-5.png}
\end{figure}
\end{frame}


\begin{frame}{}
  \begin{center}
    \fontsize{25pt}{6}\selectfont\vspace{1.2cm}
    \textcolor{subtitle}{Molecular Dynamics}
  \end{center}
\end{frame}

\begin{frame}{Molecular Dynamics}
There are three main ingredients in molecular dynamics (MD):
\begin{itemize}
\setlength\itemsep{1em}
  \item The algorithm to integrate the equations of motion,
  \item The model describing the interparticle interactions,
  \item The calculation of forces and energies from the model.
\end{itemize}
\end{frame}

\begin{frame}{Verlet Algorithm}
\begin{itemize}
\setlength\itemsep{1em}
  \item The simplest approach to obtain a numerical integration scheme is to use a Taylor series:
  \begin{equation}
  \label{eq:v-plus}
    \br_i(t+\Delta t) \approx \br_i(t) + \Delta t \mathbf{v}_i(t) + \frac{\Delta t^2}{2m_i}\mathbf{F}_i(t),
  \end{equation}
  and
  \begin{equation}
  \label{eq:v-minus}
    \br_i(t-\Delta t) \approx \br_i(t) - \Delta t \mathbf{v}_i(t) + \frac{\Delta t^2}{2m_i}\mathbf{F}_i(t).
  \end{equation}

  \item Adding Eq.~\ref{eq:v-plus} to Eq.~\ref{eq:v-minus}, we get a velocity-independent Verlet algorithm:
  \begin{equation}
    \br_i(t+\Delta t) = 2\br_i(t) - \br_i(t-\Delta t) + \frac{\Delta t^2}{m_i}\mathbf{F}_i(t).
  \end{equation}
\end{itemize}
\end{frame}

\begin{frame}{Velocity Verlet Algorithm}
\begin{itemize}
\setlength\itemsep{1em}
  \item Consider a shift in time in comparision to Eq.~\ref{eq:v-plus}:
  \begin{equation}
    \br_i(t) \approx \br_i(t+\Delta t) - \Delta t \mathbf{v}_i(t+\Delta t) + \frac{\Delta t^2}{2m_i}\mathbf{F}_i(t+\Delta t),
  \end{equation}
  which after some rearrangements takes the following form:
  \begin{equation}
    \mathbf{v}_i(t+\Delta t) = \mathbf{v}_i(t) + \frac{\Delta t}{2m_i}\left[ \mathbf{F}_i(t) + \mathbf{F}_i(t+\Delta t) \right].
  \end{equation}

  \item This scheme allows to use both evolution of the positions and velocities.
\end{itemize}
\end{frame}

\begin{frame}{Initial Conditions}
\begin{itemize}
\setlength\itemsep{1em}
  \item Initial coordinates.

  \item Initial velocities are sampled from a Maxwell-Boltzmann distribution, taking care to ensure that the sampled velocities are consistent with any constraints imposed on the system:
  \begin{equation}
    f(v) = \frac{m}{2\pi k_\mathrm{B}T}^{1/2} \e^{-mv^2 / 2 k_\mathrm{B}T},
  \end{equation}
  where $v$ is the velocity of a particle of mass $m$ at temperature $T$.

  \item An example of a Gaussian probability distribution.
\end{itemize}
\end{frame}

\begin{frame}{Potential Energy Model}
\begin{itemize}
\setlength\itemsep{1em}
  \item Commonly used model of potential energy in MD for biological macromolecules is the following:
  \begin{align}
    U(\br)
      &= \sum_\mathrm{bond} \frac{1}{2} K_\mathrm{bond} (r-r_0)^2 \\
      &+ \sum_\mathrm{angle} \frac{1}{2} K_\mathrm{bends} (\theta-\theta_0)^2 \\
      &+ \sum_\mathrm{dihedral} \sum_{n=0}^6 A_n [1+\cos(C_n \phi + \delta_n)] \\
      &+ \sum_{i,j\in \mathrm{non-bonded}} \left( 4\epsilon_{ij} \left[ \frac{\sigma_{ij}}{r_{ij}}^{12} - \frac{\sigma_{ij}}{r_{ij}}^6 \right] + \frac{q_i q_j}{r_{ij}} \right).
  \end{align}

  \item Force is given by $\mathbf{F} = -\partial{U}/\partial{\br}$.

  \item Called \textit{force field} in MD.
\end{itemize}
\end{frame}

\begin{frame}{Potential Energy Model}
\begin{itemize}
  \item Force field provides constraints into the systems.
  \begin{figure}
    \includegraphics[width=\textwidth]{fig/ff.png}
    \caption{Force field contributions.}
  \end{figure}

  \item Many force fields are available: Amber, CHARMM, Gromos.
\end{itemize}
\end{frame}

\begin{frame}{General Scheme in MD}
\begin{figure}
  \includegraphics[width=0.7\textwidth]{fig/md-scheme.png}
\end{figure}
\end{frame}

\begin{frame}[fragile]{Particle on the Wolfe-Quapp Potential}
\begin{lstlisting}
for i in range(0, args.n_samples-1):
  position[i+1] = position[i]
                + velocity[i]*dt
                + 0.5*(force/m)*dt**2
  p, f = U.eval(position[i+1])
  velocity[i+1] = velocity[i]
                + 0.5/m*dt*(force+f)
  p_energy[i+1] = p
  force = f
\end{lstlisting}
\end{frame}

\begin{frame}[fragile]{Particle on the Wolfe-Quapp Potential}
\begin{lstlisting}
python MD-WolfeQuapp.py --n_samples 10000 --dt 0.01 --mass 1
\end{lstlisting}
\vspace*{-0.3cm}
\begin{figure}
  \includegraphics[width=\textwidth]{../code/fig/md-wq.png}
\end{figure}
\end{frame}

\begin{frame}[fragile]{Particle on the M\"uller-Brown Potential}
\begin{lstlisting}
python MD-MuellerBrown.py --n_samples 10000 --dt 0.01 --mass 1
\end{lstlisting}
\vspace*{-0.3cm}
\begin{figure}
  \includegraphics[width=\textwidth]{../code/fig/md-mb.png}
\end{figure}
\end{frame}

\begin{frame}{Particle on the M\"uller-Brown Potential}
\begin{figure}
  \includegraphics[width=\textwidth]{../code/fig/md-mb-trajx.png}
\end{figure}
\end{frame}

\begin{frame}{Particle on the M\"uller-Brown Potential}
\begin{figure}
  \includegraphics[width=\textwidth]{../code/fig/md-mb-trajy.png}
\end{figure}
\end{frame}

\begin{frame}{Alanine Dipeptide}
\begin{itemize}
\setlength\itemsep{1em}
  \item Benchmark system for MD methods.
  \begin{figure}
    \includegraphics[width=0.7\textwidth]{fig/alanine.png}
    \caption{Alanine dipeptide with the dihedral angles.}
  \end{figure}

  \item Scatter plot of the dihedral angles in called the Ramachandran plot.
\end{itemize}
\end{frame}

\begin{frame}[fragile]{Alanine Dipeptide}
\begin{itemize}
\setlength\itemsep{1em}
  \item An example of the \texttt{.gro} file with alanine dipeptide.
\end{itemize}
\begin{lstlisting}
     22
    1ACE   HH31    1   1.075   1.111   1.585
    1ACE    CH3    2   1.129   1.065   1.502
    1ACE   HH32    3   1.070   0.991   1.448
    1ACE   HH33    4   1.215   1.011   1.540
    1ACE      C    5   1.162   1.167   1.395
    1ACE      O    6   1.140   1.286   1.416
    2ALA      N    7   1.230   1.118   1.289
    2ALA      H    8   1.248   1.019   1.293
    2ALA     CA    9   1.260   1.194   1.165
    2ALA     HA   10   1.304   1.125   1.093
    2ALA     CB   11   1.126   1.224   1.093
    2ALA    HB1   12   1.078   1.312   1.135
    2ALA    HB2   13   1.147   1.247   0.989
    2ALA    HB3   14   1.073   1.129   1.094
    2ALA      C   15   1.357   1.317   1.171
    2ALA      O   16   1.392   1.372   1.068
    3NME      N   17   1.404   1.353   1.294
    3NME      H   18   1.350   1.322   1.373
    3NME    CH3   19   1.517   1.445   1.315
    3NME   HH31   20   1.495   1.544   1.278
    3NME   HH32   21   1.546   1.457   1.420
    3NME   HH33   22   1.606   1.405   1.267
\end{lstlisting}
\end{frame}

\begin{frame}{Alanine Dipeptide}
\begin{itemize}
\setlength\itemsep{1em}

  \item Free energy can be calculated if the trajectory diplays ergodicity as:
  \begin{equation}
    F(\Phi,\Psi)=-\beta^{-1}\log P(\Phi,\Psi),
  \end{equation}
  where $\beta$ is the inverse temperature and $P(\Phi,\Psi)$ is the probability of occurence in a given state described by the dihedral angles.
\end{itemize}
\end{frame}

\begin{frame}{Alanine Dipeptide}
\begin{itemize}
\setlength\itemsep{1em}
  \item Difference between the ergodic and non-ergodic trajectories.
  \begin{figure}
    \includegraphics[width=\textwidth]{fig/ala1-time.png}
    \caption{Time evolution of alanine dipeptide in vacuum.}
  \end{figure}

  \item In the left, the transitions between two energy basins are \textit{rare} or \textit{infrequent}.

  \item In the right, the trajectory jumps between the states many times.
\end{itemize}
\end{frame}

\begin{frame}{Proteins in MD}
\begin{itemize}
\setlength\itemsep{1em}
  \item Proteins, DNA, viruses, membranes can be studied using MD.
  \begin{figure}
    \includegraphics[width=0.8\textwidth]{fig/pes-protein.jpg}
    \caption{Protein in the phase space.}
  \end{figure}
\end{itemize}
\end{frame}

\begin{frame}{MD Engines}
  \begin{figure}
    \includegraphics[width=0.4\textwidth]{fig/gmx.png}
  \end{figure}
\begin{itemize}
\setlength\itemsep{1em}
  \item Gromacs \url{https://www.gromacs.org/}
  \item NAMD \url{https://www.ks.uiuc.edu/Research/namd/}
  \item OpenMM \url{https://openmm.org/}
  \item LAMMPS \url{https://lammps.sandia.gov/}
\end{itemize}
\end{frame}

\begin{frame}{LAMMPS}
  \begin{quote}
    LAMMPS is a classical molecular dynamics (MD) code that models ensembles of particles in a liquid, solid, or gaseous state. It can model atomic, polymeric, biological, solid-state (metals, ceramics, oxides), granular, coarse-grained, or macroscopic systems using a variety of interatomic potentials (force fields) and boundary conditions. It can model 2d or 3d systems with only a few particles up to millions or billions.\fref{Source: LAMMPS Page: \url{https://lammps.sandia.gov/}}
  \end{quote}
\end{frame}

\begin{frame}[fragile]{LAMMPS: Executables for Linux}
  \begin{itemize}
    \item Pre-built Ubuntu Linux executable (\texttt{lammps-stable}):
    \begin{lstlisting}
  $ sudo add-apt-repository ppa:gladky-anton/lammps
  $ sudo apt-get update
  $ sudo apt-get install lammps-stable
    \end{lstlisting}

    \item Pre-built OpenSuse Linux executable (\texttt{lmp}):
    \begin{lstlisting}
  $ zypper install lammps
    \end{lstlisting}

    \item Archlinux build-script (\texttt{lmp}):
    \begin{lstlisting}
  $ git clone https://aur.archlinux.org/lammps.git
  $ cd lammps
  $ makepkg -s
  $ makepkg -i
    \end{lstlisting}

    \item Also Conda:
    \begin{lstlisting}
  $ conda config --add channels conda-forge
  $ conda create -n lammps-env
  $ conda activate lammps-env
  $ conda install lammps
    \end{lstlisting}
  \end{itemize}
  \frefw{Source: \url{https://lammps.sandia.gov/doc/Install.html}}
\end{frame}

\begin{frame}[fragile]{LAMMPS: \texttt{cmake}}
  \begin{itemize}
    \item However, it is suggested to build LAMMPS yourself.
    \item You can do this executing \texttt{./build-lammps.sh}.
  \end{itemize}
    \begin{lstlisting}
#!/bin/bash
set -e

LAMMPS_VERSION=3Mar2020
Make_CPUs=4
URL=https://github.com/lammps/lammps/archive
wget ${URL}/stable_${LAMMPS_VERSION}.zip
unzip stable_${LAMMPS_VERSION}.zip
rm -f stable_${LAMMPS_VERSION}.zip
cd lammps-stable_${LAMMPS_VERSION}

mkdir build; cd build
cmake -DPKG_MANYBODY=yes \
      -DPKG_KSPACE=yes \
      -DPKG_MOLECULE=yes \
      -DPKG_RIGID=yes \
	     ../cmake
make VERBOSE=1 -j ${Make_CPUs}
    \end{lstlisting}
\end{frame}

\begin{frame}[fragile]{LAMMPS}
  \begin{itemize}
  \setlength\itemsep{1em}
    \item The tarball for this project can be downloaded here: \href{https://ves-code.github.io/doc-v2.7-irtg-school-2020/user-doc/html/tutorial-resources/irtg-school-mainz-2020-metad.tar.gz}{click}.

    \item Contains the input file for the unbiased simulation.

    \begin{lstlisting}
  irtg-school-mainz-2020/ $ tree
  .
  |-- data.nacl
  |-- in.nacl
  |-- nacl.psf
  |-- run.sh
  |-- start.lmp
  |-- state.vmd
    \end{lstlisting}
  \end{itemize}
  \frefw{Source: \url{https://ves-code.github.io/doc-v2.7-irtg-school-2020/user-doc/html/ves_tutorials.html}}
\end{frame}

\begin{frame}{LAMMPS}
  \begin{itemize}
    \setlength\itemsep{1em}
    \item The files comes from the TRR 146 IRTG School 2020 tutorial, and were prepared by \href{https://www2.mpip-mainz.mpg.de/~valsson/}{Omar Valsson}.
    \item We consider NaCl in aqueous solution.
    \item The dissociation barrier is expected to be around 2.5 $\kT$.
    \item Collective solvent motions play an important role in the transition.
    \item The system contains 1 Na, 1 Cl, and 106 water molecules (total 320 atoms).
  \end{itemize}
\end{frame}

\begin{frame}{LAMMPS}
  \begin{itemize}
    \setlength\itemsep{1em}
    \item The system:
  \end{itemize}
  \begin{figure}
    \includegraphics[width=0.6\textwidth]{fig/NaCl.png}
    \caption{NaCl in water.}
  \end{figure}
\end{frame}

\begin{frame}[fragile]{LAMMPS}
  \begin{itemize}
    \setlength\itemsep{1em}
    \item Download VMD (\href{https://www.ks.uiuc.edu/Development/Download/download.cgi?PackageName=VMD}{click})

    \item Load the LAMMPS data file (\texttt{data.nacl}).

    \item Open VMD $\rightarrow$ Extensions $\rightarrow$ TK console.
    \begin{lstlisting}
  $ topo readlammpsdata data.nacl
  $ play state.vmd
    \end{lstlisting}

    \item We will run the simulation for 1 ns (see \texttt{start.lmp})

    \item You can do this by executing \texttt{run.sh}.
    \begin{lstlisting}
  #!/bin/bash

  lmpexec=$HOME/lammps-stable_3Mar2020/build/lmp
  ncor=2

  mpirun -np ${ncor} ${lmpexec} < start.lmp > out.lmp
    \end{lstlisting}
  \end{itemize}
\end{frame}

\begin{frame}{}
  \begin{center}
    \vspace{1.2cm}
    \includegraphics[width=0.8\textwidth]{fig/plumed-logo.png}
  \end{center}
  \frefw{Source: \url{https://www.plumed.org/}}
\end{frame}

\begin{frame}{PLUMED}
  PLUMED is an open-source, community-developed library that provides a wide range of different methods, which include:
  \begin{itemize}
    \item enhanced-sampling algorithms,
    \item free-energy methods,
    \item tools to analyze the vast amounts of data produced by molecular dynamics (MD) simulations.
  \end{itemize}
  These techniques can be used in combination with a large toolbox of collective variables that describe complex processes in physics, chemistry, material science, and biology.\fref{Source: \url{https://www.plumed.org/}}
\end{frame}

\begin{frame}[fragile]{PLUMED: Build for Linux}
  \begin{itemize}
    \item Run \texttt{build-plumed.sh} (\texttt{plumed}):
    \begin{lstlisting}
  #!/bin/bash
  set -e

  PLUMED_VERSION=2.6.2
  Make_CPUs=2
  URL=https://github.com/plumed/plumed2/archive
  wget ${URL}/v${PLUMED_VERSION}.zip
  unzip v${PLUMED_VERSION}.zip
  cd plumed2-${PLUMED_VERSION}
  INSTALL_DIR=${PWD}/install

  ./configure --prefix=${INSTALL_DIR}
  make -j ${Make_CPUs}
  make install
    \end{lstlisting}

    \item Or use Conda:
    \begin{lstlisting}
  $ conda config --add channels conda-forge
  $ conda create -n lammps-env
  $ conda activate lammps-env
  $ conda install -c conda-forge plumed
    \end{lstlisting}
  \end{itemize}
  \frefw{Source: \url{https://www.plumed.org/doc-v2.7/user-doc/html/_installation.html}}
\end{frame}

\begin{frame}[fragile]{PLUMED}
  \begin{itemize}
    \setlength\itemsep{1em}
    \item After you have finished the simulation, you should have \texttt{out.dcd} which stores the 1-ns trajectory of our system.
    \item Now we will use PLUMED to analyze our simulation.
    \item The distance between Na and Cl can be calculated using the following PLUMED action:
    \begin{lstlisting}
  distance: DISTANCE ATOMS=319,320
    \end{lstlisting}
    \item And the coordination number which represents the collective motion of the solvent wrt Na can be computed as:
    \begin{lstlisting}
  COORDINATION ...
  GROUPA=319
  GROUPB=1-318:3
  SWITCH={RATIONAL R_0=0.315 D_MAX=0.5 NN=12 MM=24}
  NLIST
  NL_CUTOFF=0.55
  NL_STRIDE=10
  LABEL=coord
  ... COORDINATION
    \end{lstlisting}
  \end{itemize}
\end{frame}

\begin{frame}[fragile]{PLUMED}
  \begin{itemize}
    \setlength\itemsep{1em}
    \item The coordination number calculates the following rational switching function:
    \begin{equation}
      f(x)=\frac{1 - \left( \frac{x}{r_0} \right)^n}{1 - \left( \frac{x}{r_0} \right)^m},
    \end{equation}
    where $r_0=0.315$, $n=12$, and $m=24$. You can see the function \href{https://www.desmos.com/calculator/nqujdsvj0z}{here.}

    \item Use this command to calculate the defined variables along the trajectory (stored in \texttt{run-driver.sh}):
    \begin{lstlisting}
  plumed driver --plumed plumed-driver.dat \
                --timestep 0.002 \
                --trajectory-stride 100 \
                --mf_dcd out.dcd
    \end{lstlisting}
  \end{itemize}
\end{frame}

\begin{frame}[fragile]{PLUMED}
  \begin{itemize}
    \setlength\itemsep{1em}
    \item You should obtain \texttt{COLVAR-driver} and see something like:
    \begin{lstlisting}
  #! FIELDS time distance coord
   0.200000 0.568067 5.506808
   0.400000 0.500148 4.994588
   0.600000 0.449778 4.931140
   0.800000 0.528272 5.105816
   1.000000 0.474371 5.089863
   1.200000 0.430620 5.091551
   1.400000 0.470374 4.993886
   1.600000 0.458768 4.940097
   1.800000 0.471886 4.952868
   2.000000 0.489058 4.897593
  .
  .
    \end{lstlisting}
  \end{itemize}
\end{frame}

\begin{frame}{LAMMPS}
  \textbf{Exercises}\\[0.2cm]
  \begin{enumerate}
    \setlength\itemsep{1em}
    \item Load the trajectory: in terminal \texttt{vmd out.dcd -psf nacl.psf}, then in VMD \texttt{pbc wrap -compound res -all}.

    \item Change the visualization of atoms and render the system.

    \item How the distance between Na and Cl changes in time?

    \item Plot distances as a histogram (normalized). What distance value is the most probable?

    \item From the histogram calculate $F=-\frac{1}{\beta}\log N$, where $N$ is the histrogram.

    \item Where the barrier is located?

    \item Perform the same analysis (from step 2) for the coordination number.
  \end{enumerate}
\end{frame}


\begin{frame}{}
  \begin{center}
    \fontsize{25pt}{6}\selectfont\vspace{1.2cm}
    \textcolor{subtitle}{Enhanced Sampling}
  \end{center}
\end{frame}

\begin{frame}{Enhanced Sampling: Source}
  \begin{itemize}
  \setlength\itemsep{1em}
    \item O. Valsson, P. Tiwary, and M. Parrinello, \textit{Enhancing Important Fluctuations: Rare Events and Metadynamics from a Conceptual Viewpoint}, Annual Review of Physical Chemistry, 67:159-184, 2016.
  \end{itemize}
\end{frame}

\begin{frame}{Long Timescales}
\begin{itemize}
\setlength\itemsep{1em}
  \item Modeling the long-timescale behavior of complex dynamical systems is a fundamental task in physical sciences.

  \item MD simulations allow us to probe the spatiotemporal details of molecular processes, but the so-called sampling problem severely limits their usefulness in practice.

  \item This sampling problem comes from the fact that a typical free energy landscape is characterized by many metastable states separated by free energy barriers that are much higher than the thermal energy $\kT$.
\end{itemize}
\end{frame}

\begin{frame}{$\kT$}
  \begin{itemize}
  \setlength\itemsep{1em}
    \item $\kT$ is used in physics as a scale factor for energy values in molecular-scale systems.
    \item Inverse of $\kT$ is called the inverse temperature and generally as $\beta=\frac{1}{\kT}$.
    \item For a system in equilibrium in canonical ensemble, the probability of the system being in state with energy $U$ is proportional to $\e^{-\beta U}$.
    \item $\kT$ is the amount of heat required to increase the thermodynamic entropy of a system by $k_{\mathrm{B}}$.
  \end{itemize}
\end{frame}

\begin{frame}{$\kT$}
  \begin{figure}
    \includegraphics[width=0.5\textwidth]{fig/kt.png}
    \caption{$\kT$ in different units.\fref{Source: \url{https://en.wikipedia.org/wiki/KT_(energy)}}}
  \end{figure}
\end{frame}

\begin{frame}{Metastable States}
  \begin{figure}
    \includegraphics[width=\textwidth]{fig/metastate.jpg}
    \caption{Slow transitions between metastable states.}
  \end{figure}
\end{frame}

\begin{frame}{Theory of Enhanced Sampling}
\begin{itemize}
\setlength\itemsep{1em}
  \item Consider a molecular system, described by microscopic coordinates $\mathbf{R}$ and a potential energy function $U(\mathbf{R})$.

  \item We limit our theory to the canonical ensemble (NVT).

  \item At equilibrium, the microscopic coordinates follow the Boltzmann distribution:
  \begin{equation}
    P(\mathbf{R}) = \e^{-\beta U(\mathbf{R})}/\int \d\mathbf{R} \,\e^{-\beta U(\mathbf{R})},
  \label{eq:pr}
  \end{equation}
  where $\beta$ is the inverse of the thermal energy.

  \item We identify a small set of coarse-grained order parameters called \textit{collective variables} (CVs):
  \begin{equation}
    \mathbf{s}(\mathbf{R}) = [s_1(\mathbf{R}), s_2(\mathbf{R}), \ldots, s_d(\mathbf{R})].
  \end{equation}
\end{itemize}
\end{frame}

\begin{frame}{Collective Variables}
  \begin{itemize}
  \setlength\itemsep{1em}
    \item Also called \textit{order parameters} or \textit{coarse-grained variables} in some publications.
    \item Should be of lower dimension than microscopic coordinates $\bR$.
    \item Can be any function of microscopic coordinates $s(\bR)$.
    \item For instance: angles, distances, contact map, dihedral angles, index along a predefined reaction coordinate, position, and many others.
    \item Can be periodic.
    \item Any differentiable function (to calculate forces) that discriminates between metastable states.
    \item Usually selecting CVs is a very difficult task.
    \item Nowadays machine learning can be used to find CVs: dimensionality reduction methods based on neural networks.
  \end{itemize}
\end{frame}

\begin{frame}{Theory of Enhanced Sampling}
\begin{itemize}
\setlength\itemsep{1em}
  \item The equilibrium marginal distribution of CVs is given by integrating out all other degrees of freedom:
  \begin{align}
  \label{eq:ps}
    P(\mathbf{s}) =
      \int \d\mathbf{R} \, \delta
        \left[
          \mathbf{s} - \mathbf{s}(\mathbf{R})
        \right]
      P(\mathbf{R}),
  \end{align}
  where we use the sifting property of the Dirac $\delta$-function.

  \item Eq.~\ref{eq:ps} is equivalent to:
  \begin{equation}
    P(\bs) = \big< \delta[\mathbf{s}-\mathbf{s(R)}] \big>,
  \end{equation}
  where $\left<\cdot\right>$ denotes an ensemble average.

  \item Up to an unimportant constant, the \textit{free energy} surface (FES) is given by:
  \begin{equation}
    F(\mathbf{s})= -\beta^{-1} \log P(\mathbf{s}).
  \end{equation}
\end{itemize}
\end{frame}

\begin{frame}{Theory of Enhanced Sampling}
\begin{itemize}
\setlength\itemsep{1em}
  \item FES contains all the relevant information about the system, such as the location and relative stability of the metastable states and the free energy barriers separating them.

  \item The relative free energy difference between two metastable states $A$ and $B$ can be determined from $F(\bs)$ by integrating the probabilities of the two states:
  \begin{equation}
    \Delta F_{A,B} = -\frac{1}{\beta}\log\frac{\int_A \d\bs\e^{-\beta F(\bs)}}{\int_B \d\bs\e^{-\beta F(\bs)}},
  \end{equation}
  where the integrations are performed over the two regions in CV space that define the two metastable states $A$ and $B$, respectively.
\end{itemize}
\end{frame}

\begin{frame}{Theory of Enhanced Sampling}
\begin{itemize}
\setlength\itemsep{1em}
  \item For an ergodic system, the FES can equivalently be computed as:
  \begin{equation}
    F(\bs) = -\frac{1}{\beta} \lim_{t\rightarrow\infty} N(\bs, t),
  \end{equation}
  where the normalized histogram is given as:
  \begin{equation}
    N(\bs,t) = \frac{\int_0^t \d t'\; \delta[\bs-\bs(\bR(t'))]}{\int_0^t \d t'}.
  \end{equation}
\end{itemize}
\end{frame}

\begin{frame}{Theory of Enhanced Sampling}
\begin{itemize}
\setlength\itemsep{1em}
  \item CV-based enhanced sampling methods overcome the sampling problem by introducing an external bias potential $V(\bs(\bR))$ acting in CV space.

  \item This leads to sampling according to a \textit{biased} distribution:
  \begin{equation}
    P_{V}(\mathbf{R}) = \frac{\e^{-\beta \left[U(\mathbf{R})+V(\mathbf{s}(\mathbf{R})) \right]}}{\int \d\mathbf{R} \,\e^{-\beta \left[U(\mathbf{R})+V(\mathbf{s}(\mathbf{R})) \right]}}
  \end{equation}

  \item Compare it to Eq.~\ref{eq:pr}.

  \item Idea of Torrie and Valleau, 1977.
\end{itemize}
\end{frame}

\begin{frame}{Theory of Enhanced Sampling}
\begin{itemize}
\setlength\itemsep{1em}
  \item At convergence, CVs follow a biased distribution:
  \begin{align}
      P_{V}(\mathbf{s}) &= \int \d\mathbf{R} \, \delta
    \left[
      \mathbf{s} - \mathbf{s}(\mathbf{R})
    \right]
  P_{V}(\mathbf{R}) \nonumber\\ &=
  \frac{ \e^{ -\beta\left[ F(\mathbf{s}) + V(\mathbf{s}) \right] } }
    {\int\d\mathbf{s} \, \e^{ -\beta\left[ F(\mathbf{s}) + V(\mathbf{s}) \right] } },
  \end{align}
  which is easier to sample.

  \item The standard reweighting from importance sampling applies here:
  \begin{equation}
    P(\bR) \propto P_V(\bR) w(\bs(\bR)) = P_V(\bR) \e^{\beta V(\bs(\bR))}.
  \end{equation}

  \item CV-based methods differ in how they construct the bias potential and which kind of biased CV sampling they obtain at convergence.
\end{itemize}
\end{frame}

\begin{frame}{Metadynamics}
\begin{itemize}
\setlength\itemsep{1em}
  \item Introduced by Laio and Parrinello in 2002.\fref{See: \url{https://doi.org/10.1073/pnas.202427399}}
  \item A dynamical scheme in which energy basins are filled in using a time-dependent potential that depends on the history of the system's trajectory.
  \item Once a basin is filled in, the system is driven into the next basin.
\end{itemize}
\begin{figure}
  \includegraphics[width=0.7\textwidth]{fig/meta.png}
\end{figure}
\end{frame}

\begin{frame}{Metadynamics}
\begin{itemize}
\setlength\itemsep{1em}
  \item Let's see how the unbiased probability distribution looks:
  \begin{equation}
    P(\bs) = \left< \prod_{\alpha=1}^n \delta\left[ f_\alpha(\br) - \bs_\alpha \right] \right>.
  \end{equation}
  \item We can replace the phase space average with a time average over a trajectory:
  \begin{equation}
    P(\bs) = \lim_{\tau\rightarrow\infty} \frac{1}{\tau} \int_0^\tau \d t\; \prod_{\alpha=1}^n \delta\left[ f_\alpha(\br(t)) - \bs_\alpha \right],
  \end{equation}
  under the assumption that the dynamics is ergodic.
\end{itemize}
\end{frame}

\begin{frame}{Metadynamics}
\begin{itemize}
\setlength\itemsep{1em}
  \item In metadynamics, we express the Dirac function as the limit of a Gaussian function:
  \begin{equation}
    \delta[x-a] = \lim_{\sigma\rightarrow 0} \frac{1}{\sqrt{2\pi\sigma^2}} \e^{-(x-a)^2/2\sigma^2}.
  \end{equation}
  \item Then the probability can be written as:
  \begin{align}
    P(\bs) &= \lim_{\tau\rightarrow\infty}\lim_{s\rightarrow 0} \frac{1}{\sqrt{2\pi\Delta s^2}\tau} \nonumber\\
           &\cdot \int_0^\tau \d t\; \prod_{\alpha=1}^n \exp\left[ -\frac{(\bs_\alpha - f_\alpha(\br(t)))^2}{2\Delta s^2} \right].
  \label{eq:metap}
  \end{align}
\end{itemize}
\end{frame}

\begin{frame}{Metadynamics}
\begin{itemize}
\setlength\itemsep{1em}
  \item Numerically, Eq.~\ref{eq:metap} is written as a discrete sum:
  \begin{align}
    P(\bs) &\approx \frac{1}{\sqrt{2\pi\Delta s^2}\tau} \nonumber\\
           &\cdot \sum_{k=0}^{N-1} \exp\left[ -\sum_{\alpha=1}^n \frac{(\bs_\alpha - f_\alpha(\br(t)))^2}{2\Delta s^2} \right].
    \label{eq:metasum}
  \end{align}
  \item Eq.~\ref{eq:metasum} suggests a form of the bias potential.
  \item Consider a bias potential of the form:
  \begin{equation}
    V(\br,t) = w \sum_{t=k\tau} \exp\left[ -\sum_{\alpha=1}^n \frac{(f_\alpha(\br) - f_\alpha(\br_\tau(t)))^2}{2\Delta s^2} \right],
    \label{eq:metav}
  \end{equation}
  where $k=1,2,\dots$ and $\tau$ is the time stride of depositing the bias.
  \item $\br_\tau(t)$ is the biased time evolution of the complete microscopic coordinates up to time $t$.
\end{itemize}
\end{frame}

\begin{frame}{Metadynamics}
\begin{itemize}
\setlength\itemsep{1em}
  \item In the simplest variant of metadynamics, the negative of Eq.~\ref{eq:metav} is taken as an approximation to the free energy:
  \begin{equation}
    F(\mathbf{q}) = -\lim_{t\rightarrow\infty} V(\mathbf{q}(\br), t).
  \end{equation}
\end{itemize}
\begin{figure}
  \includegraphics[width=0.5\textwidth]{fig/metabias.jpg}
\end{figure}
\end{frame}

\begin{frame}{Metadynamics}
\begin{itemize}
\setlength\itemsep{1em}
  \item Mapping a high-dimensional PES to a low-dimensional FES.
\end{itemize}
\begin{figure}
\centering
\begin{subfigure}{.5\textwidth}
  \centering
  \includegraphics[width=0.75\linewidth]{fig/pes}
  \caption{PES (C. Dellago)}
\end{subfigure}%
\begin{subfigure}{.5\textwidth}
  \centering
  \includegraphics[width=\linewidth]{fig/fes}
  \caption{FES (O. Valsson)}
\end{subfigure}
\end{figure}
\end{frame}

\begin{frame}{Well-Tempered Metadynamics}
\begin{itemize}
\setlength\itemsep{1em}
  \item We can build a bias potential using the following expression:
  \begin{equation}
    V(\bs)=-\left( 1-\frac{1}{\gamma} \right) F(\bs),
  \label{eq:wtmetad}
  \end{equation}
  where $\gamma\geq 1$ is a parameter called bias factor.

  \item Then we can insert Eq.~\ref{eq:wtmetad} to the formula for the biased probability distribution:
  \begin{align}
    P_V(\bs) & \propto \exp\left[ -\beta(F(\bs)+V(\bs)) \right] \nonumber\\
             & = \exp\left[ -\beta(F(\bs)-(1-1/\gamma)F(\bs)) \right] \nonumber\\
             & = \exp\left[ -\frac{\beta}{\gamma} F(\bs) \right] = [P(\bs)]^{1/\gamma},
  \label{eq:pv}
  \end{align}
  which is the so-called well-tempered distribution.
\end{itemize}
\end{frame}

\begin{frame}{Well-Tempered Metadynamics}
\begin{itemize}
\setlength\itemsep{1em}
  \item Fluctuations enhanced.
  \item Easier to cross barriers.
  \item Flatten the sampling comparing to the unbiased distribution $P(\bs)$.
\end{itemize}
\begin{figure}
  \includegraphics[width=\textwidth]{fig/enh}
\end{figure}
\end{frame}

\begin{frame}{Well-Tempered Metadynamics}
\begin{itemize}
\setlength\itemsep{1em}
  \item Using Eq.~\ref{eq:pv} and:
  \begin{equation}
    F_V(\bs) = -\frac{1}{\beta}\log P_V(\bs),
  \end{equation}
  we can obtain an expression for the free energy under the bias potential:
  \begin{equation}
    F_V(\bs) = \frac{1}{\gamma} F(\bs),
  \end{equation}
  where we ignore unimportant additive constants.

  \item We sample on an effective FES where the barriers are reduced by a factor of $\gamma$.

  \item Optimal $\gamma$ is such that effective barriers are around $\kT$.
\end{itemize}
\end{frame}

\begin{frame}{Well-Tempered Metadynamics}
\begin{itemize}
\setlength\itemsep{1em}
  \item Taking the limit $\gamma\to\infty$ gives:
  \begin{equation}
    V(\bs) = -F(\bs)
  \end{equation}
  and
  \begin{equation}
    P_V(\bs) \propto 1,
  \end{equation}
  which gives uniform sampling of CVs.

  \item Completely reduces of free energy barriers.

  \item Not optimal: Spend a lot of time sampling irrelevant regions in free energy.

  \item Better to enhance CVs with a finite value of $\gamma$.
\end{itemize}
\end{frame}

\begin{frame}{Well-Tempered Metadynamics}
\begin{itemize}
\setlength\itemsep{1em}
  \item To achieve such sampling, we can iteratively build a bias by periodically updating it:
  \begin{align}
    V_n(\bs) &= V_{n-1}(\bs) \nonumber\\
             &+ G(\bs,\bs_n)\exp\left[ -\frac{1}{\gamma-1}\beta V_{n-1}(\bs_n) \right],
  \end{align}
  where:
  \begin{equation}
    G(\bs,\bs_n) = W_0 \exp\left[ -\sum_i^d \frac{(s_i-s_{n,i})^2}{2\sigma^2_i} \right]
  \end{equation}
  is a Gaussian biasing kernel of height $W_0$ that is centered at the current CV location $\bs_n$.
\end{itemize}
\end{frame}

\begin{frame}{Well-Tempered Metadynamics}
\begin{itemize}
\setlength\itemsep{1em}
  \item The above update is performed every $N_G$ steps.

  \item Between steps $n$ and $n+1$, the bias is:
  \begin{equation}
    V_n(\bs) = \sum_{k=1}^n W_k \exp\left[ -\sum_i \frac{(s_i-s_{k,i})^2}{2\sigma^2_i} \right],
  \end{equation}
  where $W_k$ is the Gaussian height for the Gaussian deposited at step $k$:
  \begin{equation}
    W_k=W_0\exp\left[ -\frac{1}{\gamma-1} \beta V_{k-1}(\bs_k) \right],
  \end{equation}
  where the scaling factor decreases to zero as $1/n$:
  \begin{equation}
    \exp\left[ -\frac{1}{\gamma-1} \beta V_{n-1}(\bs_n) \right] \to \frac{1}{n}.
  \end{equation}
\end{itemize}
\end{frame}

\begin{frame}{Well-Tempered Metadynamics}
\begin{itemize}
\setlength\itemsep{1em}
  \item The bias reaches a quasi-stationary state, which is important for reweighting.

  \item It can be proven that in the long-time limit the bias potential is:
  \begin{equation}
    V(\bs,t)=-\left( 1-\frac{1}{\gamma} \right) F(\bs) + c(t),
  \end{equation}
  where $c(t)$ is a time-dependent constant that can be estimated during a simulation.

  \item To perform reweighting, we need to account for time-dependence of the bias:
  \begin{equation}
    P(\bR,t) = P_V(\bR)\e^{\beta [V(\bs(\bR),t)-c(t)]}.
  \end{equation}
\end{itemize}
\end{frame}

\begin{frame}{Well-Tempered Metadynamics}
\begin{itemize}
\setlength\itemsep{1em}
  \item Any thermodynamic quantity $O$ can be calculated as:
  \begin{equation}
    \left< O(\bR) \right> = \left< O(\bR) \e^{\beta [V(\bs(\bR))-c(t)]} \right>_V,
  \end{equation}
  where the time-dependent constant $c(t)$ is calculated as:
  \begin{equation}
    c(t) = \frac{1}{\beta}\log\frac{\int\d\bs\;\exp\left[ \frac{\gamma}{\gamma-1}\beta V(\bs,t) \right]}{\int\d\bs\;\exp\left[ \frac{1}{\gamma-1}\beta V(\bs,t) \right]},
  \end{equation}
  which is valid when the bias has reached a quasi-stationary state.
\end{itemize}
\end{frame}

\begin{frame}{Well-Tempered Metadynamics}
\begin{itemize}
\setlength\itemsep{1em}
  \item It can be shown that a time-independent estimator for $F(\bs)$ is given by:
  \begin{align}
    F(\bs) &= -\left( \frac{\gamma}{\gamma-1} \right)V(\bs, t) \nonumber\\
           &+ \frac{1}{\beta}\log\int\d\bs\; \exp\left[ \frac{\gamma}{\gamma-1} \beta V(\bs,t) \right]
  \end{align}

  \item The nontempered limit $\gamma\to\infty$ corresponds to:
  \begin{equation}
    V(\bs) = -F(\bs),
  \end{equation}
  which means that the bias fully compensates the underlying free energy landscape.

  \item This was the idea that guided the first version of metadynamics in which the strengh of the biasing kernel was kept constant at each iteration.
\end{itemize}
\end{frame}

\begin{frame}{Well-Tempered Metadynamics}
\begin{itemize}
\setlength\itemsep{1em}
  \item Alanine dipeptide in vacuum with backbone dihedral angles $\Phi$ and $\Psi$ as CVs.

  \item Two metastable basins (or three), $C7_{\mathrm{eq}}$ and $C7_{\mathrm{ax}}$.

  \item Barrier aroung 34 kJ/mol (16 $\kT$).

  \item Mean transition time around 28 $\mu$s.
\end{itemize}
\begin{figure}
  \includegraphics[width=\textwidth]{fig/wtala}
\end{figure}
\end{frame}

\begin{frame}{Well-Tempered Metadynamics}
\begin{itemize}
\setlength\itemsep{1em}
  \item 300 times shorter trajectory.
\end{itemize}
\begin{figure}
  \includegraphics[width=\textwidth]{fig/wtala1}
\end{figure}
\end{frame}

\begin{frame}{Well-Tempered Metadynamics}
\begin{figure}
  \includegraphics[width=0.7\textwidth]{fig/metaus}
\end{figure}
\end{frame}

\begin{frame}{Biasing Na-Cl}
\begin{itemize}
\setlength\itemsep{1em}
  \item Updated files for the enhanced sampling simulation are provided \href{https://github.com/jakryd/0800-fizobl/tree/main/project-2}{here}.
  \item We want to construct a bias potential $V(\bs)$ on the Na-Cl distance using WT-MetaD.
  \item We will run the simulation for 10 ns.
  \item The bias is given by:
  \begin{equation}
    V(\bs)=-\left(1-\frac{1}{\gamma}\right)F(\bs),
  \end{equation}
  so the system evolve under an effective free energy $F(\bs)/\gamma$.
\end{itemize}
\end{frame}

\begin{frame}[fragile]{Biasing Na-Cl}
\begin{itemize}
\setlength\itemsep{1em}
  \item Biasing the distance using PLUMED:
\end{itemize}
\begin{lstlisting}
d1: DISTANCE ATOMS=319,320
metad: METAD ...
   ARG=d1
   SIGMA=0.02
   HEIGHT=1
   BIASFACTOR=5
   TEMP=300.0
   PACE=500
   GRID_MIN=0.2
   GRID_MAX=1.0
   GRID_BIN=300
   CALC_RCT
...
\end{lstlisting}
\end{frame}

\begin{frame}[fragile]{Biasing Na-Cl}
\begin{itemize}
\setlength\itemsep{1em}
  \item We also limit the exploration of the CV space by using an upper wall bias $V_w(\bs)$:
  \begin{equation}
    V_w(s)=\kappa(s-s_0)^2
  \end{equation}
  if $s>s_0$ and 0 otherwise.
\end{itemize}
\begin{lstlisting}
d1: DISTANCE ATOMS=319,320
uwall: UPPER_WALLS ...
   ARG=d1
   AT=0.6
   KAPPA=2000.0
   EXP=2
   EPS=1
   OFFSET=0.
...
\end{lstlisting}
\end{frame}

\begin{frame}[fragile]{Biasing Na-Cl}
\begin{itemize}
\setlength\itemsep{1em}
  \item Run the simulation using: \texttt{./run.sh} (as before).
  \item Estimate the free energy profile using \texttt{HILLS} file that resulted from the simulation:
  \begin{lstlisting}
plumed sum_hills --hills HILLS --min 0.1 --max 0.8 \
                 --bin 300 --stride 100
  \end{lstlisting}
\end{itemize}
\end{frame}

\begin{frame}[fragile]{Biasing Na-Cl}
\begin{itemize}
\setlength\itemsep{1em}
  \item We should observe two states: associated and dissociated.

  \item Is the simulation converged?

  \item What energy should the system have while in the dissociated state?

  \item Why the free energy starts to rise for distance $>0.6$ nm?
\end{itemize}
\end{frame}

\begin{frame}[fragile]{Reweighting Na-Cl simulation}
\begin{itemize}
\setlength\itemsep{1em}
  \item The weights to reweight the simulation are of the following form:
  \begin{equation}
    w \propto \e^{\beta [V(\bs,t)-c(t)+V_w(\bs)]},
  \end{equation}
  where $V_w(\bs)$ is the wall bias appended to the simulation.
\end{itemize}

\begin{lstlisting}
# Read COLVAR file
distance: READ FILE=COLVAR IGNORE_TIME VALUES=d1
metad: READ FILE=COLVAR IGNORE_TIME VALUES=metad.rbias
uwall: READ FILE=COLVAR IGNORE_TIME VALUES=uwall.bias

# Define weights
weights: REWEIGHT_BIAS TEMP=300 ARG=metad.rbias,uwall.bias
\end{lstlisting}

\begin{itemize}
\setlength\itemsep{1em}
  \item A script to reweight the simulation is \texttt{plumed-driver.dat}.
  \item Execute it using: \texttt{plumed -driver --plumed plumed-driver.dat --noatoms}
\end{itemize}
\end{frame}

\begin{frame}[fragile]{Biasing Na-Cl}
\begin{itemize}
\setlength\itemsep{1em}
  \item Plot the free energy from \texttt{histo.dat}.

  \item Compare it to the free energy generated without reweighting.

  \item Does the reweighting fix the free energy?
\end{itemize}
\end{frame}

\end{document}
